\documentclass{article}
\usepackage[margin=1in]{geometry} 
\usepackage{amsmath,amsthm,amssymb}

\title{Kalman Filtering \& Sensor Fusion}
\author{UBC Rocket (TVR, Ulysses)}
\date{\today}

\begin{document}

\maketitle % Displays the title, author, and date

\section{Introduction}

The goal is to take the datastream provided from the IMU and through some process determine the current pitch and roll of Ulysses. 
We have data from the IMU's accelerometer and gyroscope, which we can refer to as follows:
$$
\text{Acceleration} =
\begin{pmatrix}
a_x\\
a_y\\
a_z
\end{pmatrix}
\ \ \ \
\text{Angular Velocity} =
\begin{pmatrix}
\omega_x\\
\omega_y\\
\omega_z
\end{pmatrix}
$$
Which we want to use to find $\theta_x, \theta_y$, the angle made between Ulysses and the x/y axes.

\section{Math}
\subsection{Notation}
Denote $\theta_{x|n,k}$ as the prediction of $\theta_x$ for time $n$ made at time $k$. For simplicity, when analysing a single dimension, we may use the notation $\theta_{n,k}$.
Similarly, we define $a_{x|t}$ and $\omega_{x|t}$ as the value of their respective measurements at time $t$.
Our goal, then, is to find $\theta_{t,t}$ at our current timestep $t$, making use of $\theta_{t,t-1}$,$\theta_{t-1,t-1}$, and the measured data.
\subsection{Accelerometer}
When Ulysses is stationary or moving at a constant velocity, we expect the magnitude of its acceleration to be approximately equal to $g=9.81$. 
We can use trigonometry to find the current angles.
\[
\begin{gathered}
\text{Roll}=\theta_x=\text{arctan}2(-a_y, -a_z)
\\\text{Pitch}=\theta_y=\text{arctan}2(a_x, \sqrt{a_y^2 + a_z^2})
\end{gathered}
\]
These angles will not be accurate if Ulysses is experiencing acceleration from other sources. 
We can judge whether or not this is the case based on how far the magnitude of the entire acceleration vector is from being $g$.
We want to fuse this data in order to come to something more accurate regardless of experienced acceleration.

\subsection{Gyroscope}
By integrating over time with respect to angular velocity, we can find the current angles of Ulysses.
$$
\theta_{x|t,t}=\theta_{x|t-1,t-1}+\Delta t\cdot\omega_{x|t}
$$
We expect the initial values of $\theta_{0,0}$ to be 0, as otherwise have no point of reference to begin our calculations.
The key issue with the gyroscope is that it reports a noise 0.1\textdegree. Over the course of a flight, this uncertainty will compound, resulting in increasingly erroneous data eventually leading to a spectacular crash.
Thus, while the gyroscope data is accurate regardless of acceleration, its noise makes it prone to drift.
\subsection{Fusion}

\end{document}